\section{Análisis de resultados}
\label{sec:analisis}

\subsection{Preparación de datos y protocolo de evaluación}
\label{sec:preparacion}

\subsubsection{Adquisición y preprocesamiento}

Ambas series se registran con periodicidad horaria desde el 1 de mayo de 2026. El balanceador de carga
entró en operación real hacia el 26 de junio de 2026; los valores previos son nulos o corresponden a una
rampa de encendido de aproximadamente una semana, en la que el tráfico crece de forma sostenida hasta
alcanzar un nivel estable (Figura~\ref{fig:raw_ramp}). Este período se descarta y el historial efectivo de
trabajo queda acotado entre el 3 de julio de 2026, 00:00 UTC, y el 26 de septiembre de 2026, 22:00 UTC:
2.063 observaciones horarias por serie, sin valores faltantes ni ceros, con frecuencia horaria estrictamente
monótona.

\begin{figure}[htbp]
  \centering
  \includegraphics[width=0.85\textwidth]{01_raw_ramp.png}
  \caption{Serie ALB sin procesar, con la rampa de encendido del balanceador previa al 3 de julio de 2026.}
  \label{fig:raw_ramp}
\end{figure}

Sobre la serie transformada (logaritmo más uno), los contrastes de Dickey-Fuller aumentado
\parencite{dickey1979} y de Kwiatkowski-Phillips-Schmidt-Shin \parencite{kwiatkowski1992} coinciden en
rechazar la estacionariedad en niveles para ambas series (ALB: $p_{\text{ADF}}=0,810$,
$p_{\text{KPSS}}=0,010$; store-service: $p_{\text{ADF}}=0,174$, $p_{\text{KPSS}}=0,010$) y en no rechazarla
tras una primera diferencia (ALB: $p_{\text{ADF}}\approx0,000$, $p_{\text{KPSS}}=0,100$; store-service:
$p_{\text{ADF}}\approx0,000$, $p_{\text{KPSS}}=0,100$), lo que justifica el orden de diferenciación de los
modelos ARIMA/SARIMA empleados.

Entre el 17 de septiembre de 2026, 18:00 UTC, y el 22 de septiembre de 2026, 13:00 UTC (116 horas), una
reversión de despliegue redujo transitoriamente el tráfico observado (Figura~\ref{fig:intervencion}). Siguiendo
el tratamiento de eventos exógenos propuesto por \textcite{boxtiao1975}, esta ventana no se descarta ni se
trata como una observación válida del proceso: se imputa reemplazando cada hora por el valor observado en
la misma hora de la semana anterior, corregido por un factor de nivel que preserva el crecimiento
orgánico semana a semana (razón entre el promedio de las 168 horas previas a la intervención y el
promedio de las 168 horas de la semana anterior a esas). Tras la reversión, la serie store-service se
estabiliza en un nivel aproximadamente 7\,\% por debajo del previo a la intervención, atribuible a un
cambio en la cantidad de destinos de balanceo activos; este desplazamiento se trata como parte del proceso
genuino y no se corrige.

\begin{figure}[htbp]
  \centering
  \includegraphics[width=0.85\textwidth]{01_intervention_zoom.png}
  \caption{Detalle de la intervención de despliegue (17 al 22 de septiembre de 2026) en ambas series.}
  \label{fig:intervencion}
\end{figure}

\subsubsection{Variables utilizadas}

Los modelos de Machine Learning se alimentan con rezagos de 1 a 24, 48, 168 y 336 horas, estadísticos
móviles (media y desvío) en ventanas de 24 y 168 horas, y variables de calendario (hora local, día de la
semana, fin de semana, feriado argentino, indicador de intervención). La correlación cruzada entre ambas
series, calculada sobre la serie transformada (logaritmo más uno y primera diferencia), alcanza su máximo
en el rezago cero ($r=0,852$; Apéndice~A), sin evidencia de una relación de adelanto-atraso aprovechable;
por esa razón, ninguna serie se utiliza como variable exógena de la otra. La estacionalidad semanal, con
medias por día de la semana que difieren menos del 3\,\% respecto de la media general en ambas series
(Figura~\ref{fig:weekly}), se considera despreciable frente a la estacionalidad diaria (fuerza estacional
de 0,771 para ALB y 0,846 para store-service, según la descomposición MSTL) y se excluye de los modelos
Prophet, NeuralProphet y del período estacional utilizado por el MASE ($m=24$).

\begin{figure}[htbp]
  \centering
  \includegraphics[width=0.8\textwidth]{01_weekly_level.png}
  \caption{Nivel medio por día de la semana, ambas series, historial efectivo.}
  \label{fig:weekly}
\end{figure}

\subsubsection{Esquema de validación y horizonte de prueba}

A diferencia de una validación cruzada expandida de múltiples pliegues, este trabajo adopta una única
ventana de validación de 48 horas (Figura~\ref{fig:split}), decisión justificada por el tiempo total de
cómputo que exige comparar 32 configuraciones por serie. La Tabla~\ref{tab:protocolo} resume las ventanas
del protocolo.

\begin{figure}[htbp]
  \centering
  \includegraphics[width=0.92\textwidth]{02_split_train_val_test.png}
  \caption{Esquema de partición temporal en entrenamiento, validación y prueba.}
  \label{fig:split}
\end{figure}

\begin{table}[htbp]
  \centering
  \caption{Ventanas del protocolo de evaluación (todos los horarios en UTC).}
  \label{tab:protocolo}
  \small
  \begin{tabular}{lll}
    \toprule
    Segmento & Rango & Extensión \\
    \midrule
    Historial efectivo & 03/07/2026 00:00 -- 26/09/2026 22:00 & 2.063 h \\
    Ajuste seguro (tuning / \textit{early stopping}) & 168 h previas a la ventana de validación & 168 h \\
    Validación (selección del modelo) & 48 h previas al inicio de la intervención & 48 h \\
    Intervención (imputada) & 17/09/2026 18:00 -- 22/09/2026 13:00 & 116 h \\
    Prueba & 22/09/2026 14:00 -- 26/09/2026 22:00 & 105 h \\
    Pronóstico final & 26/09/2026 23:00 -- 28/09/2026 22:00 & 48 h \\
    \bottomrule
  \end{tabular}
\end{table}

Para evitar la fuga de información señalada en la sección de Marco Teórico, el ajuste de hiperparámetros
de los modelos de Machine Learning y la parada anticipada de las redes profundas se realizan exclusivamente
sobre la ventana de 168 horas previa a la ventana de validación, nunca sobre esta última. El Código~\ref{cod:early_stopping}
ilustra esta separación para las arquitecturas de Deep Learning, junto con la restitución de los pesos de
la mejor época tras la parada anticipada --una corrección aplicada durante el desarrollo del trabajo, que
se documenta más adelante, al tratar los modelos de Deep Learning--.

\begin{lstlisting}[caption={Ajuste con ventana segura de fuga y restitución de la mejor época (\texttt{tsa\_final/dl\_models.py}).}, label=cod:early_stopping]
fit_part = train.iloc[:-ES_WINDOW_HOURS]      # 168 h reservadas
val_part = train.iloc[-(INPUT_CHUNK_LENGTH + ES_WINDOW_HOURS):]

tracker = _BestEpochTracker()
model = self._build(
    n_epochs=MAX_EPOCHS,
    callbacks=[EarlyStopping(monitor="val_loss", patience=PATIENCE, mode="min"), tracker],
)
model.fit(series=fit_ts, val_series=val_ts, verbose=False)
if tracker.best_state is not None:
    model.model.load_state_dict(tracker.best_state)  # no se usan los pesos post-paciencia
\end{lstlisting}

\subsection{Resultados por familia de modelos}

\subsubsection{Referencias ingenuas y modelos estadísticos clásicos}

Las cuatro referencias ingenuas y los cuatro modelos clásicos (Holt-Winters, MSTL con AutoARIMA sobre la
tendencia, AutoARIMA y el SARIMA refit del Trabajo Práctico N.\textsuperscript{o} 1) constituyen el primer
nivel de comparación. En validación, la referencia estacional ingenua ($m=24$) ya resulta muy difícil de
superar: MASE de 0,489 en ALB y de 0,480 en store-service. El AutoARIMA la supera levemente en ALB (0,475)
y queda apenas por encima en store-service (0,532). El SARIMA refit del trabajo anterior, en cambio,
obtiene un MASE de validación considerablemente peor en ambas series (0,923 y 0,647, respectivamente): su
orden fue identificado sobre datos de un mes y una rampa de tráfico distinta, y no se reoptimiza en este
trabajo. Esa desventaja en validación se revierte por completo en el conjunto de prueba, como se discute
más adelante, en la sección de selección del modelo.

\subsubsection{Machine Learning: árboles con boosting y ensambles}
\label{sec:ml}

Con hiperparámetros por defecto, ningún modelo de Machine Learning supera a la referencia estacional
ingenua en validación: el mejor de ellos en ALB (CatBoost, MASE 0,598) y en store-service (Ridge, MASE
1,201) quedan por encima de 0,489 y 0,480, respectivamente. Solo tras ajustar hiperparámetros con Optuna
--sobre la misma ventana de 48 horas luego utilizada para clasificar los modelos-- tres configuraciones de
ALB (CatBoost, Random Forest y LightGBM ajustados) logran un MASE de validación inferior a la referencia
ingenua (0,353, 0,416 y 0,421, respectivamente); en store-service, ninguna configuración ajustada alcanza
ese umbral (mínimo de 0,573 con LightGBM ajustado). Estos cinco modelos se marcan como \emph{contaminados}
en este trabajo, dado que su puntaje de validación es optimista por construcción: fueron optimizados sobre
la misma ventana que luego los clasifica. El efecto es visible en el conjunto de prueba: CatBoost ajustado,
el mejor de ALB en validación, cae al puesto 17 de 32 en prueba (MASE 2,313), y Random Forest ajustado
resulta, dentro de esta familia, el de peor desempeño relativo en prueba dentro de la serie ALB (MASE 3,455,
frente a 0,416 en validación). El \emph{stacking} de los tres mejores modelos ajustados, por su parte,
resulta el peor modelo de toda la familia de Machine Learning en ambas series y en ambos conjuntos (MASE de
prueba de 4,989 en ALB y 1,981 en store-service), lo que sugiere que combinar modelos ya sobreajustados a
la ventana de validación no mitiga, sino que propaga, ese sobreajuste. La Figura~\ref{fig:shap} muestra que
los rezagos de 1, 2, 23, 24 y 48 horas dominan la importancia de las variables según SHAP, muy por delante
de los rezagos semanales (168 horas), lo que es consistente con la estacionalidad diaria dominante
documentada en la sección de preparación de datos.

\begin{figure}[htbp]
  \centering
  \includegraphics[width=0.85\textwidth]{03_shap_alb.png}
  \caption{Importancia de variables (SHAP) del mejor modelo de Machine Learning, serie ALB.}
  \label{fig:shap}
\end{figure}

\subsubsection{Deep Learning}
\label{sec:dl}

El desarrollo de las seis arquitecturas de Deep Learning reveló un defecto metodológico corregido antes de
la selección final: el mecanismo de parada anticipada de la librería Darts conserva los pesos de la última
época evaluada, no los de la época de menor error de validación, mientras que el ajuste final sobre prueba
entrena exactamente la cantidad de épocas de la mejor época. La primera corrida comparaba, en consecuencia,
un modelo sobreentrenado en validación contra un modelo correctamente detenido en prueba. La corrección
--restituir explícitamente los pesos de la mejor época tras la parada anticipada, documentada en el
Código~\ref{cod:early_stopping}-- se aplicó a las 24 combinaciones de arquitectura, serie y semilla con
parada anticipada, y modificó únicamente los puntajes de validación; los de prueba permanecen sin cambios,
salvo en tres ajustes de TFT en los que el mecanismo de atención, no determinístico en GPU, produjo una
mejor época distinta entre corridas. La Tabla~\ref{tab:dl_correccion} resume el efecto de la corrección
sobre el MASE de validación promediado entre las dos semillas utilizadas.

\begin{table}[htbp]
  \centering
  \caption{MASE de validación antes y después de la corrección de pesos, por arquitectura de Deep Learning.}
  \label{tab:dl_correccion}
  \small
  \begin{tabular}{lrrrr}
    \toprule
    & \multicolumn{2}{c}{ALB} & \multicolumn{2}{c}{store-service} \\
    Arquitectura & Antes & Corregido & Antes & Corregido \\
    \midrule
    LSTM   & 2,337 & 1,909 & 1,017 & 0,728 \\
    N-BEATS & 1,119 & 0,980 & 1,444 & \textbf{0,638} \\
    N-HiTS  & 0,875 & \textbf{0,556} & 1,097 & 0,689 \\
    TCN     & 2,323 & 1,153 & 0,760 & 0,732 \\
    TFT     & 1,516 & 1,535 & 1,342 & 1,083 \\
    TiDE    & 0,998 & 0,790 & 1,211 & 0,862 \\
    \bottomrule
  \end{tabular}
  \\[2pt]
  \footnotesize Se destaca en negrita el mejor MASE corregido de cada serie (N-HiTS en ALB, N-BEATS en
  store-service).
\end{table}

Ninguna arquitectura de Deep Learning supera a la referencia estacional ingenua en validación tras la
corrección, aunque la distancia se reduce sustancialmente respecto de la primera corrida (por ejemplo,
N-HiTS pasa de 0,875 a 0,556 en ALB). En el conjunto de prueba, N-BEATS resulta el mejor modelo de Deep
Learning en ambas series (MASE de 1,744 en ALB y 1,247 en store-service), aunque no el mejor modelo del
proyecto completo (sección de selección del modelo, más adelante). El caso más notorio es el de LSTM en ALB: la mejor época,
para ambas semillas, es la primera (Tabla~\ref{tab:mejor_epoca}), un resultado reproducido de forma idéntica
antes y después de la corrección y que se interpreta como un hallazgo genuino --el modelo prácticamente no
llega a entrenarse antes de que la pérdida de validación comience a degradarse-- y no como un artefacto de
implementación; en prueba, esta arquitectura resulta el peor modelo de las 32 configuraciones evaluadas
para ALB (MASE de 7,817, último puesto). El Transformer de fusión temporal (TFT), por su parte, nunca ocupa
el primer puesto en ninguna serie ni conjunto, pese a exigir el mayor costo de entrenamiento de todo el
proyecto (sección de costo computacional, más adelante).

\begin{table}[htbp]
  \centering
  \caption{Mejor época (de un máximo de 100, paciencia 12) por arquitectura, serie y semilla.}
  \label{tab:mejor_epoca}
  \small
  \begin{tabular}{lrrr}
    \toprule
    Arquitectura & Serie & Semilla 42 & Semilla 43 \\
    \midrule
    LSTM   & ALB            & 1  & 1  \\
    N-BEATS & ALB           & 83 & 97 \\
    N-HiTS  & ALB           & 21 & 42 \\
    LSTM   & store-service  & 18 & 17 \\
    N-BEATS & store-service & 34 & 45 \\
    \bottomrule
  \end{tabular}
\end{table}

\subsubsection{Prophet, modelos híbridos, AutoML y modelo de fundación}

La Figura~\ref{fig:hibridos} presenta el MASE de validación de las ocho configuraciones de esta fase. El
ensamble de ponderación uniforme (\emph{equal-weight}, contaminado por construcción, ya que sus miembros se
eligen con el mismo MASE de validación que este puntaje mide) encabeza la clasificación de ALB (0,404) y
prácticamente iguala a la referencia ingenua en store-service (0,480126 frente a 0,480171). Entre los
modelos no contaminados, Prophet sin ajuste de hiperparámetros resulta el más competitivo de la fase en
ambas series (0,532 y 0,506, respectivamente), mientras que AutoGluon TimeSeries se aproxima en ALB
(0,499). El híbrido de descomposición MSTL seguido de LightGBM sobre la serie desestacionalizada constituye
el fracaso más marcado de esta fase: en store-service alcanza el peor MASE de validación de las ocho
configuraciones (1,848) y, en prueba, el peor resultado de las 32 configuraciones evaluadas para esa serie
(MASE de 3,100, último puesto de 32); en ALB su MASE de prueba (5,721) también es de los más altos del
proyecto (puesto 30 de 32). Este comportamiento es consistente con una deriva recursiva: al reconstruir la
componente estacional futura repitiendo los últimos 24 valores observados, cualquier error del componente
LightGBM sobre la tendencia se acumula sin corrección a lo largo del horizonte de prueba.

\begin{figure}[htbp]
  \centering
  \includegraphics[width=0.9\textwidth]{05_val_mae_ranking.png}
  \caption{MASE de validación, modelos de Prophet, híbridos, AutoML y de fundación.}
  \label{fig:hibridos}
\end{figure}

Chronos, evaluado en modalidad de inferencia directa sin ajuste alguno sobre las series propias, obtiene un
MASE de validación intermedio en ambas series (0,598 en ALB, 0,816 en store-service), un resultado
razonable para un modelo que nunca observó estas series durante su preentrenamiento. Su intervalo de
predicción, sin embargo, solo cubre los cuantiles 0,1 y 0,9 sobre los que fue entrenado: el intervalo del
95\,\% se deja explícitamente sin definir en el pronóstico final, en lugar de aproximarlo con un cuantil
para el que el modelo no fue entrenado.

\subsubsection{Comparación entre familias}

La Figura~\ref{fig:familias} resume el mejor MASE de validación logrado por cada familia de modelos, en
ambas series. Las Tablas~\ref{tab:comparacion_alb} y~\ref{tab:comparacion_store} presentan, de forma
compacta, el mejor modelo de cada familia junto con el modelo seleccionado, el peor modelo elegible, el
SARIMA refit del trabajo anterior y Chronos; las 32 configuraciones completas de cada serie se listan en el
Apéndice~A.

\begin{figure}[htbp]
  \centering
  \includegraphics[width=0.85\textwidth]{06_best_per_family.png}
  \caption{Mejor MASE de validación por familia de modelos, ambas series.}
  \label{fig:familias}
\end{figure}

\begin{table}[htbp]
  \centering
  \caption{Comparación compacta de modelos -- serie ALB (MAE en unidades de la serie; MASE adimensional).}
  \label{tab:comparacion_alb}
  \small
  \begin{tabular}{llrrr}
    \toprule
    Modelo & Familia & MAE val. & MASE val. & MASE prueba \\
    \midrule
    \textbf{auto\_arima (seleccionado)} & Clásico & 18.807,50 & \textbf{0,475} & 2,058 \\
    seasonal\_naive\_m24 & Ingenuo & 19.369,02 & 0,489 & 1,990 \\
    SARIMA (TP1 refit) & Fundacional TP1 & 36.524,57 & 0,923 & \textbf{1,526} \\
    catboost\_tuned\textsuperscript{a} & ML & 13.994,36 & 0,353 & 2,313 \\
    catboost\_direct & ML (directo) & 21.126,48 & 0,534 & 2,364 \\
    nhits & DL & 22.021,71 & 0,556 & 2,441 \\
    prophet\_default & Prophet & 21.066,96 & 0,532 & 2,065 \\
    mstl\_lightgbm & Híbrido & 28.341,14 & 0,716 & 5,721 \\
    autogluon\_timeseries & AutoML & 19.748,53 & 0,499 & 1,974 \\
    chronos\_bolt\_base & Fundación & 23.678,74 & 0,598 & 2,007 \\
    ensemble\_equal\_weight\textsuperscript{a} & Ensamble & 16.006,29 & 0,404 & 2,352 \\
    average & Peor elegible & 149.950,95 & 3,788 & 6,011 \\
    \bottomrule
  \end{tabular}
  \\[2pt]
  \footnotesize\textsuperscript{a}Puntaje de validación contaminado; no elegible para la selección.
\end{table}

\begin{table}[htbp]
  \centering
  \caption{Comparación compacta de modelos -- serie store-service.}
  \label{tab:comparacion_store}
  \small
  \begin{tabular}{llrrr}
    \toprule
    Modelo & Familia & MAE val. & MASE val. & MASE prueba \\
    \midrule
    \textbf{seasonal\_naive\_m24 (seleccionado)} & Ingenuo & 2.281,85 & \textbf{0,480} & 1,443 \\
    auto\_arima & Clásico & 2.525,88 & 0,532 & 1,408 \\
    SARIMA (TP1 refit) & Fundacional TP1 & 3.075,43 & 0,647 & 1,426 \\
    lightgbm\_tuned\textsuperscript{a} & ML & 2.723,67 & 0,573 & 1,770 \\
    ridge\_direct & ML (directo) & 2.905,15 & 0,611 & 1,840 \\
    nbeats & DL & 3.031,39 & 0,638 & 1,247 \\
    prophet\_default & Prophet & 2.405,88 & 0,506 & 1,083 \\
    neuralprophet & Prophet & 3.297,54 & 0,694 & \textbf{1,071} \\
    mstl\_lightgbm & Híbrido & 8.779,73 & 1,848 & 3,100 \\
    autots & AutoML & 2.412,97 & 0,508 & 1,596 \\
    chronos\_bolt\_base & Fundación & 3.877,73 & 0,816 & 1,673 \\
    ensemble\_equal\_weight\textsuperscript{a} & Ensamble & 2.281,64 & 0,480 & 1,612 \\
    drift & Peor elegible & 17.022,66 & 3,582 & 2,587 \\
    \bottomrule
  \end{tabular}
  \\[2pt]
  \footnotesize\textsuperscript{a}Puntaje de validación contaminado; no elegible para la selección. Se
  incluye NeuralProphet además del mejor Prophet por ser, con amplio margen, el mejor modelo de prueba de
  toda la serie (ver la sección siguiente, de selección del modelo).
\end{table}

\subsection{Selección del modelo y consistencia entre validación y prueba}
\label{sec:seleccion}

La regla de selección adoptada (Código~\ref{cod:seleccion}) recorre el conjunto de 32 modelos por serie,
descarta los marcados como contaminados y selecciona el de menor MAE de validación entre los restantes.
Se marcan como contaminados los cinco modelos ajustados con Optuna sobre la ventana de validación en cada
serie (\texttt{catboost\_tuned}, \texttt{random\_forest\_tuned} y \texttt{lightgbm\_tuned} en ALB;
\texttt{ridge\_tuned}, \texttt{lightgbm\_tuned} y \texttt{xgboost\_tuned} en store-service), el
\emph{stacking} construido a partir de esos modelos ajustados, y el ensamble de ponderación uniforme, cuyos
miembros se eligen con el propio MAE de validación. Prophet ajustado (\texttt{prophet\_tuned}) se excluye de
esta lista porque su ajuste se realiza sobre la ventana segura de 168 horas descrita en la sección de
preparación de datos, nunca sobre la ventana de validación.

\begin{lstlisting}[caption={Regla de selección: menor MAE de validación entre los modelos elegibles (\texttt{tsa\_final/selection.py}).}, label=cod:seleccion]
def select_model(pool: pd.DataFrame, series: str) -> SelectionResult:
    """El modelo seleccionado es el de menor MAE de validacion entre los
    elegibles (la regla adoptada). El conjunto de prueba nunca se usa para
    seleccionar."""
    elig = eligible_val(pool, series)
    return SelectionResult(series=series, selected=elig.iloc[0],
                            runner_up=elig.iloc[1], worst=elig.iloc[-1])
\end{lstlisting}

Bajo esta regla, el modelo seleccionado para ALB es \texttt{auto\_arima} y para store-service es la propia
referencia estacional ingenua. La comparación contra el SARIMA refit del trabajo anterior arroja una mejora
del 48,5\,\% en validación para ALB (MASE de 0,475 frente a 0,923) y del 25,8\,\% en validación para
store-service (0,480 frente a 0,647); en ambos casos, esa ventaja se revierte en el conjunto de prueba: el
modelo seleccionado resulta un 34,9\,\% peor que el SARIMA refit en ALB y un 1,2\,\% peor en store-service.
Comparado contra la referencia estacional ingenua, el modelo seleccionado de ALB mejora un 2,9\,\% en
validación pero empeora un 3,4\,\% en prueba.

La Figura~\ref{fig:consistencia} y el coeficiente de correlación de Spearman entre el orden de validación y
el orden de prueba cuantifican esta divergencia: para ALB, considerando las 32 configuraciones, $\rho=0,237$
($p=0,192$), no significativo al 5\,\%; restringido a los 27 modelos elegibles, $\rho=0,346$ ($p=0,077$),
tampoco significativo. Para store-service, en cambio, la consistencia es moderada y estadísticamente
significativa: $\rho=0,654$ ($p<0,001$) sobre las 32 configuraciones, y $\rho=0,692$ ($p<0,001$) sobre las
27 elegibles. Esta significancia agregada, sin embargo, no impide divergencias puntuales severas: NeuralProphet,
decimocuarto en validación de store-service, resulta el mejor modelo de prueba de toda la serie.

\begin{figure}[htbp]
  \centering
  \includegraphics[width=0.85\textwidth]{06_val_vs_test_scatter.png}
  \caption{Dispersión entre el orden de validación y el orden de prueba, ambas series.}
  \label{fig:consistencia}
\end{figure}

La Tabla~\ref{tab:sensibilidad} contrasta la regla adoptada con dos criterios alternativos, a título
exclusivamente informativo: el que se obtendría si se permitieran modelos contaminados, y el que
seleccionaría un oráculo que conociera de antemano el MAE de prueba --un criterio inaplicable en la
práctica, dado que el pronóstico se produce antes de observar el conjunto de prueba--. Cada criterio elige
un ganador distinto en ambas series, lo que evidencia que la elección del criterio de selección no es un
detalle menor del protocolo.

\begin{table}[htbp]
  \centering
  \caption{Análisis de sensibilidad del criterio de selección (no constituye una regla de selección alternativa).}
  \label{tab:sensibilidad}
  \small
  \begin{tabular}{lll}
    \toprule
    Criterio & Serie ALB & Serie store-service \\
    \midrule
    Regla adoptada (MAE val., elegibles) & auto\_arima & seasonal\_naive\_m24 \\
    Incluyendo modelos contaminados & catboost\_tuned & ensemble\_equal\_weight \\
    Oráculo de MAE de prueba\textsuperscript{a} & SARIMA (TP1 refit) & neuralprophet \\
    \bottomrule
  \end{tabular}
  \\[2pt]
  \footnotesize\textsuperscript{a}No utilizable al momento de pronosticar; se incluye solo como referencia
  de contexto.
\end{table}

\subsection{Costo computacional}
\label{sec:costo}

El costo de ajuste no fue el criterio de selección, pero resulta relevante para evaluar la viabilidad
operativa de cada familia. La Tabla~\ref{tab:costo} contrasta los modelos de mayor costo con el modelo
seleccionado. El Transformer de fusión temporal (TFT) exige, por lejos, el mayor tiempo de ajuste del
proyecto (entre 175 y 292 segundos por serie y conjunto, sumado sobre las dos semillas utilizadas), sin
alcanzar el primer puesto en ninguna comparación; AutoGluon TimeSeries exhibe un costo fijo similar,
producto de su límite de tiempo configurado (300 segundos por llamada), independientemente del tamaño de
la serie.

\begin{table}[htbp]
  \centering
  \caption{Tiempo de ajuste de los modelos de mayor costo computacional (segundos).}
  \label{tab:costo}
  \footnotesize
  \begin{tabular}{lp{3.6cm}rr}
    \toprule
    Modelo & Serie & Ajuste (val.) & Reajuste (prueba) \\
    \midrule
    TFT & ALB & 292,41 & 195,67 \\
    TFT & store-service & 271,62 & 175,47 \\
    AutoGluon TimeSeries & ALB & 161,81 & 140,03 \\
    AutoGluon TimeSeries & store-service & 119,65 & 147,21 \\
    auto\_arima & ALB (seleccionado) & 97,52 & 78,93 \\
    auto\_arima & store-service & 79,47 & 46,09 \\
    seasonal\_naive\_m24 & ambas (seleccionado en store-service) & $<0,001$ & $<0,001$ \\
    \bottomrule
  \end{tabular}
\end{table}

\subsection{Pronóstico final a 48 horas}

Con el modelo seleccionado por serie reajustado sobre la totalidad del historial disponible (entrenamiento
más validación, con la intervención imputada), se produce un pronóstico de 48 horas posteriores al final
del historial, entre el 26 de septiembre de 2026 a las 23:00 UTC y el 28 de septiembre de 2026 a las 22:00
UTC. La Figura~\ref{fig:overlay} contrasta, sobre el conjunto de prueba, las trayectorias del modelo
seleccionado, del peor modelo elegible, del SARIMA refit del trabajo anterior y de Chronos, contra el valor
observado; la Figura~\ref{fig:forecast_final} muestra el pronóstico de 48 horas propiamente dicho, con sus
intervalos de predicción.

\begin{figure}[htbp]
  \centering
  \includegraphics[width=0.92\textwidth]{06_test_overlay_forecast.png}
  \caption{Pronósticos sobre el conjunto de prueba: modelo seleccionado, peor modelo, SARIMA (TP1 refit) y Chronos, ambas series.}
  \label{fig:overlay}
\end{figure}

\begin{figure}[htbp]
  \centering
  \includegraphics[width=0.92\textwidth]{06_final_forecast.png}
  \caption{Pronóstico final de 48 horas con intervalos de predicción, ambas series.}
  \label{fig:forecast_final}
\end{figure}

Los intervalos de predicción provienen de fuentes distintas según el modelo: AutoARIMA y SARIMA aportan
intervalos nativos de 80\,\% y 95\,\%; la referencia estacional ingenua y el peor modelo elegible utilizan
intervalos empíricos, construidos a partir de los residuos agrupados de la ventana de validación de 48
horas, y por lo tanto constantes en cada hora del horizonte; Chronos aporta únicamente el intervalo del
80\,\%, dado que fue entrenado solo sobre los cuantiles 0,1 y 0,9, y su intervalo del 95\,\% se deja sin
definir en lugar de aproximarse con un cuantil no entrenado.

\pendiente{Serie 3 (a cargo de otro integrante del grupo). Los resultados de esta serie --comparación por
familia de modelos, selección y pronóstico final-- se incorporarán en esta sección una vez completado el
trabajo del integrante responsable.}
